\section{Team contributions}
\label{app:contributions}

This appendix documents the division of work, deliverables, and verification records for all five
team members.

\subsection{Member A (Ruxuan ZHOU): Data preparation and experimental protocol}

Member A downloaded and preprocessed the UCI raw dataset, designed the fixed experimental protocol
(random seed, stratification strategy, and leakage controls), defined the 18-variable and
24-variable feature sets with documented exclusion criteria, and generated the split manifest and
configuration files. Verification passed 23/23 delivery checks.

\noindent\textit{Deliverables:}
\begin{itemize}[leftmargin=*,noitemsep,topsep=2pt]
  \item \path{data/processed/students.csv} (4,424 rows; checksums recorded in \path{data/manifest.json})
  \item \path{data/splits/split_manifest.csv} (\texttt{row\_id}, \texttt{dataset\_split}, \texttt{cv\_fold} assignments)
  \item \path{configs/protocol.json} and \path{configs/feature_sets.json}
  \item \path{src/A_data/prepare_data.py} and \path{src/A_data/verify_A.py}
  \item \path{docs/handoffs/A.md}
\end{itemize}

\subsection{Member B (Mingxi ZHU): Model training and hyperparameter tuning}

Member B implemented training pipelines for logistic regression, decision tree, and random forest
models, ran five-fold stratified cross-validation grid search on the training split, froze all six
model pipelines plus the majority-class baseline prior to test access, and pre-specified the
24-variable random forest as the representative model for interpretation. Verification passed 16/16
delivery checks, with all 35 validation folds independently recomputed to within $10^{-12}$.

\noindent\textit{Deliverables:}
\begin{itemize}[leftmargin=*,noitemsep,topsep=2pt]
  \item \path{src/B_models/train_models.py}, \path{src/B_models/verify_B.py},\\
  \path{src/B_models/test_b_models.py}
  \item \path{models/*.joblib} (six frozen model pipelines plus majority baseline)
  \item \path{results/B_validation/cv_metrics.csv}, \path{selected_params.json},\\
  \path{model_catalog.json}
  \item \path{docs/handoffs/B.md}
\end{itemize}

\subsection{Member C (Boyi WEI): Final evaluation and interpretation}

Member C applied the frozen pipelines to the held-out test split in a one-shot evaluation, computed
macro-F1, class-level precision and recall, accuracy, and log-loss, generated all five
interpretation figures, and computed permutation importance for the pre-specified representative
random forest. Verification passed 14/14 delivery checks; the prediction table contains 6,195 rows
(885 test students $\times$ 7 pipelines), and summary metrics recompute from stored predictions
within a maximum absolute error of $8.88 \times 10^{-16}$.

\noindent\textit{Deliverables:}
\begin{itemize}[leftmargin=*,noitemsep,topsep=2pt]
  \item \path{src/C_evaluation/evaluate_models.py} and \path{src/C_evaluation/verify_C.py}
  \item \path{results/C_final/test_metrics.csv} (seven model-level summary rows evaluated on 885 test students)
  \item \path{results/C_final/test_predictions.csv} (6,195 prediction and probability rows)
  \item \path{results/C_final/figures/*.png} (five visualisation figures)
  \item \path{results/C_final/permutation_importance_representative.csv},\\
  \path{results/C_final/logistic_dropout_coefficients.csv}, and \path{results/C_final/findings.md}
  \item \path{docs/handoffs/C.md}
\end{itemize}

\subsection{Member D (Mengqi YAN): Report integration and documentation}

Member D integrated technical materials from Members A, B, and C into the final report, drafted the
project proposal and reproducibility appendix, cross-checked reported figures against Member C's
output files, compiled references, and ensured adherence to experimental reporting standards.

\noindent\textit{Deliverables:}
\begin{itemize}[leftmargin=*,noitemsep,topsep=2pt]
  \item \path{docs/report/report.pdf} (this document)
  \item \path{docs/report/main_report.tex}, \path{docs/report/code_appendix.tex},\\
  \path{docs/report/references.bib}
  \item \path{docs/proposal/proposal.pdf}
\end{itemize}

\subsection{Member E (Yushu WU): Presentation and dissemination}

Member E cross-checked the report's numerical results against Member C's evaluation outputs (including
dropout counts, precision, and recall), prepared the presentation slides incorporating the report's
key figures, authored the bilingual English--Chinese presentation script for all five members,
coordinated and completed the recorded group presentation, and designed the A0 project poster.

\noindent\textit{Deliverables:}
\begin{itemize}[leftmargin=*,noitemsep,topsep=2pt]
  \item \path{recording/Student-Dropout-Prediction.mp4}
  \item \texttt{slides/Team-Zero bias hero-PPT.pptx}
  \item \texttt{slides/Team-Zero bias hero-Poster.pptx}
  \item \texttt{slides/Team-Zero bias hero-Script.docx}
\end{itemize}

\section{Reproducibility and code appendix}
\label{app:code}

\subsection{Repository layout}

The full codebase, configurations, frozen models, evaluation outputs, and presentation materials are
hosted in a public GitHub repository:
\begin{itemize}[leftmargin=*,noitemsep,topsep=2pt]
  \item \textbf{Repository:} \url{https://github.com/GeorgeZhou888/DSC5001-Student-Dropout-Prediction}
  \item \textbf{Branch:} \texttt{main}
  \item \textbf{Data licence:} The UCI source dataset is licensed under CC BY 4.0; no separate licence is specified for the project code.
\end{itemize}

Key project files are organised by stage:
\begin{itemize}[leftmargin=*,noitemsep,topsep=2pt]
  \item Data preparation: \path{src/A_data/prepare_data.py}, \path{src/A_data/verify_A.py}.
  \item Protocol and feature definitions: \path{configs/protocol.json}, \path{configs/feature_sets.json}.
  \item Model training and checks: \path{src/B_models/train_models.py},\\
  \path{src/B_models/verify_B.py}, \path{src/B_models/test_b_models.py}.
  \item Final evaluation and checks: \path{src/C_evaluation/evaluate_models.py},\\
  \path{src/C_evaluation/verify_C.py}.
  \item Final outputs: \path{results/C_final/test_metrics.csv},\\
  \path{results/C_final/test_predictions.csv}, interpretation CSVs, and figures in \path{results/C_final/figures/}.
\end{itemize}

\subsection{Environment}

All three computational stages ran under Python 3.14.x with dependencies pinned in\\
\path{environment/requirements-lock.txt}. The recorded runs used different operating systems and
patch versions: Stage A used Python 3.14.6 on macOS ARM64, Stage B used Python 3.14.5 on Windows
AMD64, and Stage C used Python 3.14.3 on Windows AMD64. The run manifests preserve these exact
runtime details rather than treating the environments as identical.

\subsection{Commands}

All commands below are executed from the repository root and are platform-neutral apart from the
local Python executable name.

\begin{lstlisting}[language=bash,caption={Data preparation and verification}]
python src/A_data/prepare_data.py
python src/A_data/verify_A.py
python src/A_data/verify_A.py --reproduce
\end{lstlisting}

\begin{lstlisting}[language=bash,caption={Model verification and pre-evaluation checks}]
python -m unittest discover -s src/B_models -p test_b_models.py -v
python src/B_models/verify_B.py --recheck-cv
python src/C_evaluation/evaluate_models.py --check
python src/C_evaluation/verify_C.py
\end{lstlisting}

The held-out test evaluation was executed once using the frozen Stage B pipelines, with all outputs
archived under \path{results/C_final/}:

\begin{lstlisting}[language=bash,caption={One-shot held-out test evaluation}]
python src/C_evaluation/evaluate_models.py
\end{lstlisting}

By design, \path{evaluate_models.py} refuses to overwrite an existing output directory unless an
explicit override flag (\texttt{-\/-force}) is supplied, protecting the archived one-shot evaluation
results from accidental regeneration.

\subsection{Protocol summary}

\begin{table}[H]
\centering
\caption{Fixed analysis protocol.}
\label{tab:protocol}
\begin{tabular}{ll}
\toprule
Item & Setting \\
\midrule
Raw records & 4,424 \\
Train/test counts & 3,539 / 885 \\
Outer split & Stratified 80/20, seed 42 \\
Training validation & Five-fold stratified cross-validation \\
Feature sets & 18 enrolment variables; 24 enrolment plus semester-1 variables \\
Models & Logistic regression, decision tree, random forest \\
Baseline & Majority class \\
Selection metric & Macro-F1 \\
Test-set use & One final evaluation after model freezing \\
\bottomrule
\end{tabular}
\end{table}

\subsection{Provenance notes}

The final evaluation manifest (\path{results/C_final/run_manifest.json}) records the frozen Stage B
reference, input SHA-256 digests and accepted line-ending states, model hashes, output hashes,
runtime environment, and execution command. The Stage C verification report confirms 14 of 14 checks
passing: the archived prediction table holds 6,195 rows (885 test students across seven pipelines),
and every summary metric in \path{results/C_final/test_metrics.csv} recomputes directly from those
stored predictions, establishing internal consistency and full provenance for the final project
deliverables.
